\documentclass[letterpaper,11pt]{article}
\usepackage[T1]{fontenc}
\usepackage[utf8]{inputenc}
\usepackage{lmodern}
\usepackage[french]{babel}
\usepackage[margin=1in,headheight=28pt]{geometry}
\renewcommand{\familydefault}{\sfdefault}
\usepackage{amsmath,amssymb,amsthm,bm,enumitem}
\usepackage[dvipsnames]{xcolor}
\usepackage[small,compact]{titlesec}
\usepackage{fancyhdr,fancyvrb}
\usepackage[colorlinks=true,allcolors=blue]{hyperref}
\pagestyle{fancy}
\fancyhead[L]{STT-2100 \\ Régression}
\fancyhead[R]{A.-A. Vallée \\ Université Laval}
\fancyfoot[C]{\thepage}
\setlength{\parindent}{0pt}
\setlength{\parskip}{5pt}
\fvset{fontsize=\small,frame=single,framesep=2mm}
\newcommand{\esp}{\mathbb{E}}
\newcommand{\var}{\operatorname{Var}}
\newcommand{\bo}[1]{\bm{#1}}
\newcommand{\N}{\mathcal{N}}
% Définir \avecsolutions avant l'input pour une version corrigée.
\newcommand{\prof}[1]{\ifdefined\avecsolutions\par{\color{blue}\textbf{Solution :} #1}\par\fi}
\begin{document}
\begin{center}{\Large Exercices - Chapitre 11}\end{center}
\begin{enumerate}[leftmargin=*,wide,label=\bfseries Exercice \arabic*:]
\item Le modèle $Y_i=\beta_0+\beta_1x_i+\beta_2x_i^2+\varepsilon_i$ est-il un modèle de régression linéaire? Écrivez sa matrice de modèle et le test permettant de comparer au modèle de premier degré. Pourquoi conserver $x$ si $x^2$ est retenu?
\prof{Il est linéaire en ses paramètres. Les lignes de $\bo X$ sont $(1,x_i,x_i^2)$. Testez $H_0:\beta_2=0$ par un test $t$ ou $F$ à une contrainte, sur les mêmes observations. Le maintien de $x$ respecte la hiérarchie et évite de dépendre artificiellement de l'origine choisie.}
\item On ajuste $\hat y=5+2x-0{,}10x^2$. Calculez la pente locale à $x=4$ et le point où elle s'annule. Cette dernière valeur peut-elle être interprétée hors du domaine observé?
\prof{La dérivée est $2-0{,}20x$, donc $1{,}2$ à $x=4$; elle s'annule à $x=10$. C'est un maximum de la courbe ajustée, mais son interprétation exige un domaine de données et un modèle pertinents à cet endroit.}
\item Pour $Y_i=\alpha+\beta e^{-x_i/\theta}+\varepsilon_i$, indiquez pourquoi l'ajustement est non linéaire en les paramètres. Quels problèmes des valeurs initiales peuvent-elles provoquer? Pourquoi transformer la réponse n'est-il pas en général équivalent à minimiser la somme des carrés sur l'échelle originale?
\prof{Le paramètre $\theta$ intervient dans l'exponentielle; il n'existe pas de matrice de modèle fixe donnant une combinaison linéaire de tous les paramètres. Des départs inadéquats peuvent provoquer non-convergence ou minimum local. Une transformation change l'échelle des erreurs, leur variance et le critère minimisé.}
\item Trois observations indépendantes ont des variances proportionnelles à $(1,4,9)$. Quels poids relatifs utiliser? Écrivez l'estimateur des moindres carrés pondérés pour $\var(\bo\varepsilon)=\sigma^2\bo W^{-1}$ et l'estimateur de $\sigma^2$.
\prof{Les poids sont proportionnels à $(1,1/4,1/9)$. $\hat\bo\beta=(\bo X^\top\bo W\bo X)^{-1}\bo X^\top\bo W\bo Y$ et $s_W^2=\sum_iw_i e_i^2/(n-p')$. Multiplier tous les poids par une même constante ne change pas les coefficients ajustés.}
\item Dans un modèle AR(1) stationnaire, $\varepsilon_t=\rho\varepsilon_{t-1}+u_t$, avec $\rho=0{,}6$ et $\var(u_t)=4$. Calculez la variance des erreurs et leurs covariances aux retards 1 et 2. Pourquoi les erreurs-types usuelles de \texttt{lm()} ne sont-elles pas adaptées?
\prof{$\var(\varepsilon_t)=4/(1-0{,}36)=6{,}25$; les covariances valent $3{,}75$ et $2{,}25$. La covariance n'est pas $\sigma^2\bo I$, alors que les erreurs-types usuelles reposent sur cette forme. Il faut modéliser la covariance ou utiliser une méthode d'inférence appropriée.}
\item Avec la fiche R, comparez un modèle linéaire et quadratique pour \texttt{cars}, ajustez un modèle exponentiel aux données simulées, puis une régression pondérée aux données hétéroscédastiques. Pour un modèle pondéré, quels résidus tracer pour examiner la dispersion?
\prof{Utilisez \texttt{I(speed\string^2)}, \texttt{nls()} et l'argument \texttt{weights} de \texttt{lm()}. Les résidus pondérés $\sqrt{w_i}e_i$, ou les résidus standardisés du modèle pondéré, permettent de vérifier la variance sur l'échelle rendue homogène. Vérifiez également la forme de la moyenne.}
\end{enumerate}
\end{document}
